\bigskip
\noindent Autre piste de l'onglet Findings pour ce thème, dans une cote qui n'a pas de section ici :
\begin{itemize}\small\setlength{\itemsep}{0pt}
\item Cote~2 : si $i : X \to X \times \mathbf P^r$ est une section constante et $Y \subset X$ un diviseur dont la classe $x' \in A^1(X)$ est non nulle mais tuée par $(r-1)!$, alors $i_!(\mathcal O_Y)$ est une classe non nulle de $K(X \times \mathbf P^r)$ de classe de Chern totale~$1$, ce qui n'est possible que pour $r \geq 3$. Le \q{0} final et les mots qui entourent la conclusion de la page~13 sont incertains ou illisibles, et la page nomme la classe $[\mathcal O_{Y \times \mathbf P^r}]$ là où le calcul porte sur $Y \times \{\text{point}\}$. La lecture de la page est à vérifier sur le fac-similé. Mathlib n'a pas le groupe $K$ des faisceaux cohérents ni les classes de Chern. On se reportera aux notes de Grothendieck (\lot{2}{1}{9}).\piste{2-trivial-total-chern-class}
\end{itemize}
